\section{Sliding Window XOR}
\label{sec:cses-sw-xor}

\problemstatement{Given an array of $n$ integers and a window size $k$,
output the bitwise XOR of every contiguous window of size $k$.}

\begin{patternbox}
XOR is its own inverse ($x\oplus x=0$), so, like sum, the window XOR slides
in $O(1)$: XOR in the entering element and XOR out the leaving one.
Equivalently, a \textbf{prefix-XOR} array answers each window as
$\textit{pre}[i+k]\oplus\textit{pre}[i]$.
\end{patternbox}

\subsection*{Prefix XOR}

Let $\textit{pre}[0]=0$ and $\textit{pre}[i]=a[0]\oplus\cdots\oplus
a[i-1]$. Then the XOR of $a[l..r]$ is $\textit{pre}[r+1]\oplus
\textit{pre}[l]$, because the shared prefix $a[0..l-1]$ cancels itself
(each bit XORed twice returns to $0$). A window of size $k$ starting at $i$
is $\textit{pre}[i+k]\oplus\textit{pre}[i]$.

\begin{ideabox}[Self-Inverse Means Cancellation]
Because $x\oplus x=0$, overlapping prefixes cancel exactly, so a
prefix-XOR array reduces any window XOR to a single operation. The same
self-inverse property lets a running XOR drop the departing element by
simply XORing it again.
\end{ideabox}

\subsection*{Algorithm}

\begin{enumerate}
  \item Build the prefix-XOR array (or maintain a running XOR).
  \item For each window $i$, output
        $\textit{pre}[i+k]\oplus\textit{pre}[i]$.
\end{enumerate}

\subsection*{Dry run}

\dryrun{$a=[1,2,3,4]$, $k=2$.}

\begin{itemize}
  \item $\textit{pre}=[0,1,3,0,4]$.
  \item Window $[1,2]$: $\textit{pre}[2]\oplus\textit{pre}[0]=3$.
        $[2,3]$: $\textit{pre}[3]\oplus\textit{pre}[1]=0\oplus1=1$.
        $[3,4]$: $\textit{pre}[4]\oplus\textit{pre}[2]=4\oplus3=7$.
  \item Output $3,1,7$.
\end{itemize}

\subsection*{C++ implementation}

\begin{lstlisting}
#include <bits/stdc++.h>
using namespace std;

int main() {
    int n, k;
    cin >> n >> k;
    vector<int> a(n);
    for (int& x : a) cin >> x;

    vector<int> pre(n + 1, 0);
    for (int i = 0; i < n; i++) pre[i + 1] = pre[i] ^ a[i];

    for (int i = 0; i + k <= n; i++)
        cout << (pre[i + k] ^ pre[i]) << " ";     // window XOR
    cout << "\n";
    return 0;
}
\end{lstlisting}

\begin{complexitybox}
\textbf{Time Complexity.} $O(n)$. \textbf{Space Complexity.} $O(n)$ for the
prefix array (or $O(1)$ with a running XOR).
\end{complexitybox}

\begin{takeawaybox}
XOR joins sum in the ``invertible, slides in $O(1)$'' category thanks to
being its own inverse. Prefix XOR is the range-query form and reappears in
counting subarrays with a target XOR. Contrast with OR (next), which is
\emph{not} invertible and needs per-bit counts.
\end{takeawaybox}
